\chapter{2004年浙江卷（文）}
\section{单选题}
\begin{problem}
若 \(U = \{1, 2, 3, 4\}\)，\(M = \{1, 2\}\)，\(N = \{2, 3\}\)，则 \(\complement_U(M \cup N)\) =（\quad）

\choices
    {\(\{1, 2, 3\}\)}
    {\(\{4\}\)}
    {\(\{1, 3, 4\}\)}
    {\(\{2\}\)}
\end{problem}

\begin{answer}
B
\end{answer}

\begin{solution}
由题意，\(M\cup N=\{1,2,3\}\)，所以全集 \(U\) 中除去 \(\{1,2,3\}\) 后剩余 \(\{4\}\)，即
\(\complement_U(M\cup N)=\{4\}\)，故选 B.
\end{solution}

\begin{problem}
直线 \(y = 2\) 与直线 \(x + y - 2 = 0\) 的夹角是（\quad）

\choices
    {\(\frac{\pi}{4}\)}
    {\(\frac{\pi}{3}\)}
    {\(\frac{\pi}{2}\)}
    {\(\frac{3\pi}{4}\)}
\end{problem}

\begin{answer}
A
\end{answer}

\begin{solution}
直线 \(y=2\) 的斜率为 \(0\)，直线 \(x+y-2=0\) 化为 \(y=-x+2\)，斜率为 \(-1\). 设两直线的夹角为 \(\theta\)，则
\(\tan\theta=\left|\frac{-1-0}{1+(-1)\cdot0}\right|=1\). 取两直线的锐角夹角，得 \(\theta=\frac{\pi}{4}\)，故选 A.
\end{solution}

\begin{problem}
已知等差数列 \(\{a_n\}\) 的公差为 \(2\)，若 \(a_1\)，\(a_3\)，\(a_4\) 成等比数列，则 \(a_2 =\)（\quad）

\choices
    {\(-4\)}
    {\(-6\)}
    {\(-8\)}
    {\(-10\)}
\end{problem}

\begin{answer}
B
\end{answer}

\begin{solution}
设 \(a_1=x\)，则 \(a_3=x+4\)，\(a_4=x+6\). 由于 \(a_1\)，\(a_3\)，\(a_4\) 成等比数列，有
\((x+4)^2=x(x+6)\)，即 \(x^2+8x+16=x^2+6x\)，所以 \(x=-8\). 因此 \(a_2=a_1+2=-6\)，故选 B.
\end{solution}

\begin{problem}
已知向量 \(\bs{a} = (3, 4)\)，\(\bs{b} = (\sin \alpha, \cos \alpha)\)，且 \(\bs{a} \parallel \bs{b}\)，则 \(\tan \alpha =\)（\quad）

\choices
    {\(\frac{3}{4}\)}
    {\(-\frac{3}{4}\)}
    {\(\frac{4}{3}\)}
    {\(-\frac{4}{3}\)}
\end{problem}

\begin{answer}
A
\end{answer}

\begin{solution}
因为 \(\bs{a}\parallel\bs{b}\)，所以两向量的行列式为零，即
\(3\cos\alpha-4\sin\alpha=0\). 若 \(\cos\alpha=0\)，上式不成立，故 \(\cos\alpha\neq0\)，从而 \(\tan\alpha=\frac{3}{4}\)，故选 A.
\end{solution}

\begin{problem}
点 \(P\) 从 \(\left(1, 0\right)\) 出发，沿单位圆 \(x^{2} + y^{2} = 1\) 逆时针方向运动 \(\frac{2\pi}{3}\) 弧长到达 \(Q\) 点，则 \(Q\) 的坐标为（\quad）

\choices
    {\(\left(-\frac{1}{2}, \frac{\sqrt{3}}{2}\right)\)}
    {\(\left(-\frac{\sqrt{3}}{2}, -\frac{1}{2}\right)\)}
    {\(\left(-\frac{1}{2}, -\frac{\sqrt{3}}{2}\right)\)}
    {\(\left(-\frac{\sqrt{3}}{2}, \frac{1}{2}\right)\)}
\end{problem}

\begin{answer}
A
\end{answer}

\begin{solution}
单位圆的半径为 \(1\)，所以弧长 \(\frac{2\pi}{3}\) 所对的圆心角为 \(\frac{2\pi}{3}\). 点 \(P\) 的极角为 \(0\)，逆时针转过该角后，点 \(Q\) 的坐标为
\(\left(\cos\frac{2\pi}{3},\sin\frac{2\pi}{3}\right)=\left(-\frac{1}{2},\frac{\sqrt{3}}{2}\right)\)，故选 A.
\end{solution}

\begin{problem}
曲线 \(y^{2} = 4x\) 关于直线 \(x = 2\) 对称的曲线方程是（\quad）

\choices
    {\(y^{2} = 8 - 4x\)}
    {\(y^{2} = 4x - 8\)}
    {\(y^{2} = 16 - 4x\)}
    {\(y^{2} = 4x - 16\)}
\end{problem}

\begin{answer}
C
\end{answer}

\begin{solution}
设对称后的曲线上一点为 \((x,y)\)，它关于直线 \(x=2\) 的对称点为 \((4-x,y)\). 该对称点在原曲线 \(y^2=4x\) 上，所以
\(y^2=4(4-x)=16-4x\)，故选 C.
\end{solution}

\begin{problem}
若 \(\left(\sqrt{x} + \frac{2}{\sqrt[3]{x}}\right)^{n}\) 展开式中存在常数项，则 \(n\) 的值可以是（\quad）

\choices
    {\(8\)}
    {\(9\)}
    {\(10\)}
    {\(12\)}
\end{problem}

\begin{answer}
C
\end{answer}

\begin{solution}
在展开式中取 \(k\) 个 \(\frac{2}{\sqrt[3]{x}}\)，取 \(n-k\) 个 \(\sqrt{x}\)，所得项的 \(x\) 的指数为
\(\frac{n-k}{2}-\frac{k}{3}=\frac{3n-5k}{6}\). 该项为常数项当且仅当 \(3n-5k=0\)，即 \(k=\frac{3n}{5}\). 由于 \(k\) 必须是整数且满足 \(0\leqslant k\leqslant n\)，所以 \(n\) 必须是 \(5\) 的倍数. 给出的选项中只有 \(n=10\) 满足，故选 C.
\end{solution}

\begin{problem}
在 \(\triangle ABC\) 中，“\(A > 30^{\circ}\)”是“\(\sin A > \frac{1}{2}\)”的（\quad）

\choices
    {充分而不必要条件}
    {必要而不充分条件}
    {充分必要条件}
    {既不充分也不必要条件}
\end{problem}

\begin{answer}
B
\end{answer}

\begin{solution}
三角形内角满足 \(0<A<180^{\circ}\). 在此范围内，\(\sin A>\frac{1}{2}\) 等价于 \(30^{\circ}<A<150^{\circ}\)，因此必有 \(A>30^{\circ}\). 但 \(A>30^{\circ}\) 不一定推出 \(A<150^{\circ}\)，例如 \(A=160^{\circ}\) 时 \(\sin A<\frac{1}{2}\). 所以“\(A>30^{\circ}\)”是“\(\sin A>\frac{1}{2}\)”的必要而不充分条件，故选 B.
\end{solution}

\begin{problem}
若函数 \(f(x) = \log_{a}(x + 1)\)（\(a > 0\)，\(a \neq 1\)）的定义域和值域都是 \([0, 1]\)，则 \(a =\)（\quad）

\choices
    {\(\frac{1}{3}\)}
    {\(\sqrt{2}\)}
    {\(\frac{\sqrt{2}}{2}\)}
    {\(2\)}
\end{problem}

\begin{answer}
D
\end{answer}

\begin{solution}
函数的定义域被题设限定为 \([0,1]\). 当 \(a>1\) 时，\(f(x)=\log_a(x+1)\) 在 \([0,1]\) 上递增，所以值域为 \([f(0),f(1)]=[0,\log_a2]\). 要使值域为 \([0,1]\)，必须有 \(\log_a2=1\)，从而 \(a=2\). 当 \(0<a<1\) 时，\(f(1)=\log_a2<0\)，不可能得到值域 \([0,1]\)，故选 D.
\end{solution}

\begin{problem}
如图，在正三棱柱 \(ABC - A_{1}B_{1}C_{1}\) 中已知 \(AB = 1\)，\(D\) 在棱 \(BB_{1}\) 上，且 \(BD = 1\)，若 \(AD\) 与平面 \(AA_{1}C_{1}C\) 所成的角为 \(\alpha\)，则 \(\alpha =\)（\quad）

\bitmapinlinefigure[4.0cm][max width=3.1cm]{img/2004/zhejiang_liberal/q10_fig1.png}

\choices
    {\(\frac{\pi}{3}\)}
    {\(\frac{\pi}{4}\)}
    {\(\arcsin \frac{\sqrt{10}}{4}\)}
    {\(\arcsin \frac{\sqrt{6}}{4}\)}
\end{problem}

\begin{answer}
D
\end{answer}

\begin{solution}
取空间直角坐标系，使
\(A=(0,0,0)\)，\(B=(1,0,0)\)，\(C=\left(\frac{1}{2},\frac{\sqrt{3}}{2},0\right)\)，并令正三棱柱的侧棱方向为 \(z\) 轴. 由 \(BD=1\) 得 \(D=(1,0,1)\). 平面 \(AA_1C_1C\) 的一个单位法向量可取 \(\bs{n}=\left(-\frac{\sqrt{3}}{2},\frac{1}{2},0\right)\)，而 \(\overrightarrow{AD}=(1,0,1)\). 设直线 \(AD\) 与该平面的夹角为 \(\alpha\)，则
\(\sin\alpha=\frac{|\overrightarrow{AD}\cdot\bs{n}|}{|\overrightarrow{AD}|}=\frac{\frac{\sqrt{3}}{2}}{\sqrt{2}}=\frac{\sqrt{6}}{4}\). 因而 \(\alpha=\arcsin\frac{\sqrt{6}}{4}\)，故选 D.
\end{solution}

\begin{problem}
若椭圆 \(\frac{x^{2}}{a^{2}} + \frac{y^{2}}{b^{2}} = 1\)（\(a > b > 0\)）的左、右焦点分别为 \(F_1\)，\(F_2\)，线段 \(F_1F_2\) 被点 \(\left(\frac{b}{2}, 0\right)\) 分成 \(5:3\) 两段，则此椭圆的离心率为（\quad）

\choices
    {\(\frac{16}{17}\)}
    {\(\frac{4\sqrt{17}}{17}\)}
    {\(\frac{4}{5}\)}
    {\(\frac{2\sqrt{5}}{5}\)}
\end{problem}

\begin{answer}
D
\end{answer}

\begin{solution}
设椭圆的焦点为 \(F_1=(-c,0)\)，\(F_2=(c,0)\)，其中 \(c>0\). 记点 \(R=\left(\frac{b}{2},0\right)\). 由于 \(R\) 把 \(F_1F_2\) 分成较长的 \(5\) 与较短的 \(3\) 两段，有
\(\frac{c+\frac{b}{2}}{c-\frac{b}{2}}=\frac{5}{3}\). 解得 \(c=2b\). 又椭圆满足 \(a^2=b^2+c^2\)，所以 \(a^2=5b^2\)，从而
\(e=\frac{c}{a}=\frac{2b}{\sqrt{5}b}=\frac{2\sqrt{5}}{5}\)，故选 D.
\end{solution}

\begin{problem}
若 \(f(x)\) 和 \(g(x)\) 都是定义在实数集 \(\R\) 上的函数，且方程 \(x - f[g(x)] = 0\) 有实数解，则 \(g[f(x)]\) 不可能是（\quad）

\choices
    {\(x^{2} + x - \frac{1}{5}\)}
    {\(x^{2} + x + \frac{1}{5}\)}
    {\(x^{2} - \frac{1}{5}\)}
    {\(x^{2} + \frac{1}{5}\)}
\end{problem}

\begin{answer}
B
\end{answer}

\begin{solution}
由方程 \(x-f(g(x))=0\) 有实数解，存在实数 \(x_0\) 使 \(f(g(x_0))=x_0\). 令 \(y_0=g(x_0)\)，则
\(g(f(y_0))=g(x_0)=y_0\). 因此，函数 \(g(f(x))\) 必须有实数不动点.

若选项 B 成立，则不动点方程为 \(x^2+x+\frac{1}{5}=x\)，即 \(x^2+\frac{1}{5}=0\)，没有实数解. 所以选项 B 不可能，故选 B.
\end{solution}

\section{填空题}
\begin{problem}
已知 \(f(x) = \begin{cases}
1, & x \geqslant 0,\\
-1, & x < 0.
\end{cases}\) 则不等式 \(x + (x + 2)\cdot f(x + 2) \leqslant 5\) 的解集是\fillinblank{}.
\end{problem}

\begin{answer}
\(\left(-\infty,\frac{3}{2}\right]\)
\end{answer}

\begin{solution}
当 \(x+2\geqslant0\)，即 \(x\geqslant-2\) 时，\(f(x+2)=1\)，原不等式化为 \(2x+2\leqslant5\)，所以 \(-2\leqslant x\leqslant\frac{3}{2}\). 当 \(x+2<0\)，即 \(x<-2\) 时，\(f(x+2)=-1\)，原不等式化为 \(x-(x+2)=-2\leqslant5\)，所以该情形下所有 \(x<-2\) 都满足不等式. 合并两种情形，解集为 \(\left(-\infty,\frac{3}{2}\right]\).
\end{solution}

\begin{problem}
已知平面上三点 \(A\)，\(B\)，\(C\) 满足 \(\left|\overrightarrow{AB}\right| = 3\)，\(\left|\overrightarrow{BC}\right| = 4\)，\(\left|\overrightarrow{CA}\right| = 5\)，则 \(\overrightarrow{AB}\cdot\overrightarrow{BC} + \overrightarrow{BC}\cdot\overrightarrow{CA} + \overrightarrow{CA}\cdot\overrightarrow{AB}\) 的值等于\fillinblank{}.
\end{problem}

\begin{answer}
\(-25\)
\end{answer}

\begin{solution}
由三点首尾相接，有
\(\overrightarrow{AB}+\overrightarrow{BC}+\overrightarrow{CA}=\bs{0}\). 两边平方得
\(0=|\overrightarrow{AB}|^2+|\overrightarrow{BC}|^2+|\overrightarrow{CA}|^2+2\left(\overrightarrow{AB}\cdot\overrightarrow{BC}+\overrightarrow{BC}\cdot\overrightarrow{CA}+\overrightarrow{CA}\cdot\overrightarrow{AB}\right)\).
代入三条边长 \(3\)，\(4\)，\(5\)，得到
\(0=3^2+4^2+5^2+2S=50+2S\)，所以 \(S=-25\).
\end{solution}

\begin{problem}
已知平面 \(\alpha \perp \beta\)，\(\alpha \cap \beta = l\)，\(P\) 是空间一点，且 \(P\) 到 \(\alpha\)，\(\beta\) 的距离分别是 \(1\)，\(2\)，则点 \(P\) 到 \(l\) 的距离为\fillinblank{}.
\end{problem}

\begin{answer}
\(\sqrt{5}\)
\end{answer}

\begin{solution}
在过点 \(P\) 且垂直于交线 \(l\) 的平面内，平面 \(\alpha\) 与 \(\beta\) 截得两条互相垂直的直线，它们的交点为该截平面与 \(l\) 的交点. 点 \(P\) 到这两条直线的距离分别为 \(1\)，\(2\)，因此点 \(P\) 到交点的距离为
\(\sqrt{1^2+2^2}=\sqrt{5}\).
\end{solution}

\begin{problem}
设坐标平面内有一个质点从原点出发，沿 \(x\) 轴跳动. 每次向正方向或负方向跳 \(1\) 个单位，经过 \(5\) 次跳动质点落在点 \((3, 0)\)（允许重复过此点）处，则质点不同的运动方法共有\fillinblank{} 种.
\end{problem}

\begin{answer}
\(5\)
\end{answer}

\begin{solution}
设向正方向跳了 \(r\) 次，向负方向跳了 \(5-r\) 次. 终点横坐标为 \(3\)，故
\(r-(5-r)=3\)，解得 \(r=4\). 只需从 \(5\) 次跳动中选出其中一次向负方向跳，便确定一种运动方法，因此方法数为 \(\mathrm C_5^1=5\).
\end{solution}

\section{解答题}
\begin{problem}
已知数列 \(\{a_n\}\) 的前 \(n\) 项和为 \(S_n\)，\(S_n = \frac{1}{3}(a_n - 1)\)（\(n \in \N^*\)）.

\begin{enumerate}
\item 求 \(a_1\)，\(a_2\)；
\item 求证数列 \(\{a_n\}\) 是等比数列.
\end{enumerate}
\end{problem}

\begin{solution}
\begin{enumerate}
\item 取 \(n=1\)，由 \(S_1=a_1\) 得 \(a_1=\frac{1}{3}(a_1-1)\)，所以 \(a_1=-\frac{1}{2}\). 取 \(n=2\)，有 \(a_1+a_2=\frac{1}{3}(a_2-1)\)，代入 \(a_1=-\frac{1}{2}\) 得 \(a_2=\frac{1}{4}\).
\item 当 \(n\geqslant2\) 时，将相邻两式相减，得
\(S_n-S_{n-1}=a_n=\frac{1}{3}(a_n-a_{n-1})\). 因而 \(2a_n=-a_{n-1}\)，即 \(a_n=-\frac{1}{2}a_{n-1}\). 因此从 \(a_1=-\frac{1}{2}\) 出发，数列的每一项都是前一项的 \(-\frac{1}{2}\) 倍，所以数列 \(\{a_n\}\) 是等比数列，公比为 \(-\frac{1}{2}\).
\end{enumerate}
\end{solution}

\begin{problem}
在 \(\triangle ABC\) 中，角 \(A\)，\(B\)，\(C\) 所对的边分别为 \(a\)，\(b\)，\(c\)，且 \(\cos A = \frac{1}{3}\).

\begin{enumerate}
\item 求 \(\sin^{2}\frac{B + C}{2} + \cos 2A\) 的值；
\item 若 \(a = \sqrt{3}\)，求 \(bc\) 的最大值.
\end{enumerate}
\end{problem}

\begin{solution}
\begin{enumerate}
\item 因为 \(B+C=\pi-A\)，所以
\(\sin^2\frac{B+C}{2}=\sin^2\left(\frac{\pi}{2}-\frac{A}{2}\right)=\cos^2\frac{A}{2}=\frac{1+\cos A}{2}=\frac{2}{3}\). 又 \(\cos2A=2\cos^2A-1=\frac{2}{9}-1=-\frac{7}{9}\)，故
\(\sin^2\frac{B+C}{2}+\cos2A=\frac{2}{3}-\frac{7}{9}=-\frac{1}{9}\).
\item 由余弦定理，\(a^2=b^2+c^2-2bc\cos A=b^2+c^2-\frac{2}{3}bc\). 当 \(a=\sqrt3\) 时，
\(3=b^2+c^2-\frac{2}{3}bc\geqslant2bc-\frac{2}{3}bc=\frac{4}{3}bc\)，所以 \(bc\leqslant\frac{9}{4}\). 当 \(b=c=\frac{3}{2}\) 时等号成立，且满足题设三角形条件，因此 \(bc\) 的最大值为 \(\frac{9}{4}\).
\end{enumerate}
\end{solution}

\begin{problem}
如图，已知正方形 \(ABCD\) 和矩形 \(ACEF\) 所在的平面互相垂直，\(AB = \sqrt{2}\)，\(AF = 1\)，\(M\) 是线段 \(EF\) 的中点.

\begin{enumerate}
\item 求证：\(AM \parallel\) 平面 \(BDE\)；
\item 求证：\(AM \perp\) 平面 \(BDF\)；
\item 求二面角 \(A - DF - B\) 的大小.
\end{enumerate}

\bitmapfigure[max width=4.0cm]{img/2004/zhejiang_liberal/q19_fig1.png}
\end{problem}

\begin{solution}
取空间直角坐标系，使正方形所在平面为 \(z=0\)，并设
\(A=(0,0,0)\)，\(B=(\sqrt{2},0,0)\)，\(C=(\sqrt{2},\sqrt{2},0)\)，\(D=(0,\sqrt{2},0)\). 由于矩形 \(ACEF\) 与正方形所在平面垂直，且 \(AF=1\)，可取
\(E=(\sqrt{2},\sqrt{2},1)\)，\(F=(0,0,1)\)，从而 \(M=\left(\frac{\sqrt{2}}{2},\frac{\sqrt{2}}{2},1\right)\).

\begin{enumerate}
\item \(\overrightarrow{AM}=\left(\frac{\sqrt{2}}{2},\frac{\sqrt{2}}{2},1\right)\). 平面 \(BDE\) 内的两条相交直线可取 \(BD\)，\(BE\)，其方向向量为
\(\overrightarrow{BD}=(-\sqrt{2},\sqrt{2},0)\)，\(\overrightarrow{BE}=(0,\sqrt{2},1)\). 因而平面 \(BDE\) 的一个法向量为
\(\overrightarrow{BD}\times\overrightarrow{BE}=(\sqrt{2},\sqrt{2},-2)\). 计算得
\(\overrightarrow{AM}\cdot(\sqrt{2},\sqrt{2},-2)=1+1-2=0\). 又
\((A-B)\cdot(\sqrt{2},\sqrt{2},-2)=-2\neq0\)，所以 \(A\) 不在平面 \(BDE\) 内，故 \(AM\parallel\) 平面 \(BDE\).
\item 平面 \(BDF\) 内的两条相交直线可取 \(BD\)，\(BF\)，其中 \(\overrightarrow{BF}=(-\sqrt{2},0,1)\). 因而平面 \(BDF\) 的一个法向量为
\(\overrightarrow{BD}\times\overrightarrow{BF}=(\sqrt{2},\sqrt{2},2)=2\overrightarrow{AM}\). 所以 \(AM\perp\) 平面 \(BDF\).
\item 平面 \(ADF\) 的一个法向量为 \((1,0,0)\)，平面 \(BDF\) 的一个法向量为 \((\sqrt{2},\sqrt{2},2)\). 设二面角 \(A-DF-B\) 为 \(\theta\)，则
\(\cos\theta=\frac{|(1,0,0)\cdot(\sqrt{2},\sqrt{2},2)|}{\sqrt{2+2+4}}=\frac{1}{2}\). 因而 \(\theta=\frac{\pi}{3}\).
\end{enumerate}
\end{solution}

\begin{problem}
某地区有 \(5\) 个工厂，由于用电紧缺，规定每个工厂在一周内必须选择某一天停电（选哪一天是等可能的）. 假定工厂之间的选择互不影响.

\begin{enumerate}
\item 求 \(5\) 个工厂均选择星期日停电的概率；
\item 求至少有两个工厂选择同一天停电的概率.
\end{enumerate}
\end{problem}

\begin{solution}
\begin{enumerate}
\item 每个工厂有 \(7\) 种等可能选择，所有选择的总数为 \(7^5\). 其中五个工厂均选择星期日的选择只有一种，所以所求概率为 \(\frac{1}{7^5}\).
\item 先求五个工厂选择不同日期的概率. 五个工厂依次选择不同日期的方案数为 \(7\cdot6\cdot5\cdot4\cdot3\)，所以至少有两个工厂选择同一天的概率为
\(1-\frac{7\cdot6\cdot5\cdot4\cdot3}{7^5}=1-\frac{2520}{16807}=\frac{2041}{2401}\).
\end{enumerate}
\end{solution}

\begin{problem}
已知 \(a\) 为实数，\(f(x) = (x^{2} - 4)(x - a)\).

\begin{enumerate}
\item 求导数 \(f'(x)\)；
\item 若 \(f'(-1) = 0\)，求 \(f(x)\) 在 \(\left[-2, 2\right]\) 上的最大值和最小值；
\item 若 \(f(x)\) 在 \(\left(-\infty, -2\right]\) 和 \(\left[2, +\infty\right)\) 上都是递增的，求 \(a\) 的取值范围.
\end{enumerate}
\end{problem}

\begin{solution}
\begin{enumerate}
\item 展开得 \(f(x)=x^3-ax^2-4x+4a\)，所以 \(f'(x)=3x^2-2ax-4\).
\item 由 \(f'(-1)=0\)，得 \(3+2a-4=0\)，所以 \(a=\frac{1}{2}\). 此时
\(f(x)=x^3-\frac{1}{2}x^2-4x+2\)，且 \(f'(x)=3x^2-x-4=(3x-4)(x+1)\). 区间 \([-2,2]\) 内的可能极值点为 \(-2\)，\(-1\)，\(\frac{4}{3}\)，\(2\). 逐一计算得
\(f(-2)=0\)，\(f(-1)=\frac{9}{2}\)，\(f\left(\frac{4}{3}\right)=-\frac{50}{27}\)，\(f(2)=0\). 因此最大值为 \(\frac{9}{2}\)，最小值为 \(-\frac{50}{27}\).
\item 要使函数在给定两个区间上递增，需使 \(f'(x)=3x^2-2ax-4\geqslant0\) 对 \(x\leqslant-2\) 及 \(x\geqslant2\) 均成立. 令
\(h(x)=\frac{3x^2-4}{2x}=\frac{3}{2}x-\frac{2}{x}\)，则 \(h'(x)=\frac{3}{2}+\frac{2}{x^2}>0\).

当 \(x\geqslant2\) 时，\(f'(x)\geqslant0\) 等价于 \(a\leqslant h(x)\). 由于 \(h\) 在 \([2,+\infty)\) 上递增，故只需 \(a\leqslant h(2)=2\). 当 \(x\leqslant-2\) 时，除以负数 \(2x\) 后不等号反向，\(f'(x)\geqslant0\) 等价于 \(a\geqslant h(x)\). 由于 \(h\) 在 \((-\infty,-2]\) 上递增，故只需 \(a\geqslant h(-2)=-2\). 综合得 \(-2\leqslant a\leqslant2\).
\end{enumerate}
\end{solution}

\begin{problem}
已知双曲线的中心在原点，右顶点为 \(A(1, 0)\)，点 \(P\)，\(Q\) 在双曲线的右支上，点 \(M(m, 0)\) 到直线 \(AP\) 的距离为 \(1\).

\begin{enumerate}
\item 若直线 \(AP\) 的斜率为 \(k\)，且 \(\left|k\right| \in \left[\frac{\sqrt{3}}{3}, \sqrt{3}\right]\)，求实数 \(m\) 的取值范围；
\item 当 \(m = \sqrt{2} + 1\) 时，\(\triangle APQ\) 的内心恰好是点 \(M\)，求此双曲线的方程.
\end{enumerate}
\end{problem}

\begin{solution}
\begin{enumerate}
\item 直线 \(AP\) 的方程为 \(y=k(x-1)\)，即 \(kx-y-k=0\). 点 \(M(m,0)\) 到该直线的距离为
\(\frac{|k||m-1|}{\sqrt{k^2+1}}=1\)，所以
\(|m-1|=\frac{\sqrt{k^2+1}}{|k|}=\sqrt{1+\frac{1}{k^2}}\).
令 \(t=|k|\)，则 \(t\in\left[\frac{\sqrt{3}}{3},\sqrt{3}\right]\)，且 \(\sqrt{1+t^{-2}}\) 随 \(t\) 增大而减小. 因此
\(\frac{2\sqrt{3}}{3}\leqslant|m-1|\leqslant2\). 从而
\(m\in\left[-1,1-\frac{2\sqrt{3}}{3}\right]\cup\left[1+\frac{2\sqrt{3}}{3},3\right]\).
\item 设双曲线方程为 \(x^2-\frac{y^2}{b^2}=1\)，其中 \(b>0\). 因为 \(M\) 是三角形 \(APQ\) 的内心，\(AM\) 是顶点 \(A\) 的角平分线. 又 \(A\)，\(M\) 均在 \(x\) 轴上，所以 \(AP\)，\(AQ\) 关于 \(x\) 轴对称. 双曲线也关于 \(x\) 轴对称，故可设
\(P=(x_0,y_0)\)，\(Q=(x_0,-y_0)\)，其中 \(x_0>1\)，\(y_0>0\).

\(PQ\) 为直线 \(x=x_0\). 由于 \(M\) 是三角形 \(APQ\) 的内点，且 \(A=(1,0)\)，故 \(1<m<x_0\). 内心到三边距离相等，点 \(M\) 到直线 \(AP\) 的距离为 \(1\)，所以 \(x_0-m=1\)，从而 \(x_0=m+1=2+\sqrt{2}\). 设直线 \(AP\) 的斜率为正数 \(k\)，则
\(\frac{k(m-1)}{\sqrt{k^2+1}}=1\). 代入 \(m-1=\sqrt{2}\)，得 \(\frac{\sqrt{2}k}{\sqrt{k^2+1}}=1\)，从而 \(k^2=1\). 由于 \(k>0\)，故 \(k=1\)，进而 \(y_0=k(x_0-1)=1+\sqrt{2}\).

把点 \(P=(2+\sqrt{2},1+\sqrt{2})\) 代入双曲线，得
\(b^2=\frac{(1+\sqrt{2})^2}{(2+\sqrt{2})^2-1}=\frac{3+2\sqrt{2}}{5+4\sqrt{2}}=\frac{1+2\sqrt{2}}{7}=\frac{1}{2\sqrt{2}-1}\). 因此双曲线方程为
\(x^2-(2\sqrt{2}-1)y^2=1\). 此时 \(x_0>1\)，且 \(y_0>0\)，故 \(P\)，\(Q\) 确在该双曲线的右支上；由对称性，点 \(M\) 到 \(AP\)，\(AQ\)，\(PQ\) 的距离均为 \(1\)，条件相容.
\end{enumerate}
\end{solution}
